% Eksempelfil for å se Høstlys på LaTeX. Kan ikke kompileres alene.
\documentclass[a4paper,11pt]{article}
\usepackage{siunitx}
\usepackage{notatstil}

\begin{document}

\section{Banen til en planet}

Keplers 2.\ lov sier at arealet planeten sveiper over per tidsenhet er det samme i hele banen. I de to vendepunktene står farten $v$ normalt på avstanden $r$ til sola, og vi får \emph{arealhastigheten} direkte, se likning~\eqref{eq:areal-vendepunkt}.
\begin{equation}
    \frac{1}{2}\, r_0 v_0 \Delta t = \frac{1}{2}\, r_1 v_1 \Delta t
    \qquad\Longrightarrow\qquad
    v_1 = v_0\,\frac{r_0}{r_1} .
    \label{eq:areal-vendepunkt}
\end{equation}

Energibevaring mellom de to vendepunktene gir
\begin{align*}
    \frac{1}{2} m v_0^2 - \frac{\gamma M m}{r_0} &= \frac{1}{2} m v_1^2 - \frac{\gamma M m}{r_1} , \\
    \frac{1}{2} v_0^2 \left(1 - \frac{r_0^2}{r_1^2}\right) &= \gamma M \left(\frac{1}{r_0} - \frac{1}{r_1}\right) \quad \text{etter at vi har satt inn $v_1$.}
\end{align*}

Gravitasjonskraften fra sola på planeten er $\vv{F}_1 = -\frac{\gamma M m}{r^2}\,\hat{\symbf{r}}$, og ved jordoverflata er tyngdeakselerasjonen \qty{9,81}{m/s^2}. Vi bruker tre steg:
\begin{itemize}
    \item finn uttrykket for $v_1$ fra arealloven,
    \item sett det inn i energilikningen,
    \item løs for $v_0$ og sett inn tall helt til slutt.
\end{itemize}

\end{document}
